\section{Joint Deep-Chain Geometry: Exact Spectra and Depth Dependence}
\label{app:rigid-chain}
\subsection{Joint group action and the normal operator}
\begin{definition}
For $L\geq2$,
\begin{align}
W_1&\in\R_*^{r\times d},&
W_i&\in{\rm GL}(r)\ (1<i<L),&
W_L&\in\R_*^{m\times r},\\
\pi_L(W)&=W_L\cdots W_1,&
\mathcal G_L&={\rm GL}(r)^{L-1},\\
W_1'&=Q_1^{-1}W_1,&
W_i'&=Q_i^{-1}W_iQ_{i-1},&
W_L'&=W_LQ_{L-1}.
\end{align}
Let $\Omega=(\Omega_1,\ldots,\Omega_{L-1})$ and
\begin{align}
(\mathcal V_W\Omega)_1&=-\Omega_1W_1,\\
(\mathcal V_W\Omega)_i&=W_i\Omega_{i-1}-\Omega_iW_i,\quad1<i<L,\\
(\mathcal V_W\Omega)_L&=W_L\Omega_{L-1},\\
(\mathcal V_W^*u)_i&=W_{i+1}^\top u_{i+1}-u_iW_i^\top .
\label{eq:rigid-chain-vertical}
\end{align}
\end{definition}

\begin{proposition}[Joint projection, not pairwise projection]
\label{prop:joint-chain}
\begin{align}
D\pi_L(W)[u]&=\sum_{i=1}^L L_i u_i R_i,\\
L_i&=W_L\cdots W_{i+1},&
R_i&=W_{i-1}\cdots W_1,\\
L_L&=I_m,&R_1&=I_d,\\
\ker D\pi_L(W)&=\qrange\mathcal V_W,\\
\Omega_u&=(\mathcal V_W^*\mathcal V_W)^{-1}\mathcal V_W^*u,&
\Pi_Wu&=u-\mathcal V_W\Omega_u .
\end{align}
For $\mathcal N_W=\mathcal V_W^*\mathcal V_W$,
\begin{align}
(\mathcal N_W\Omega)_i
&=W_{i+1}^\top W_{i+1}\Omega_i+\Omega_iW_iW_i^\top\notag\\
&\quad-W_i\Omega_{i-1}W_i^\top
-W_{i+1}^\top\Omega_{i+1}W_{i+1}.
\label{eq:rigid-block-tridiagonal}
\end{align}
Terms involving $\Omega_0,\Omega_L$ are omitted.
\end{proposition}
\begin{proof}
\begin{align}
D\pi_L[\mathcal V_W\Omega]
&=-L_1\Omega_1W_1
+\sum_{i=2}^{L-1}L_i(W_i\Omega_{i-1}-\Omega_iW_i)R_i
+W_L\Omega_{L-1}R_L=0,\\
\sum_i\ii{u_i}{(\mathcal V_W\Omega)_i}
&=\sum_{i=1}^{L-1}
\ii{W_{i+1}^\top u_{i+1}-u_iW_i^\top}{\Omega_i}.
\end{align}
If $\mathcal V_W\Omega=0$,
\begin{align}
\Omega_1W_1=0&\Longrightarrow\Omega_1=0,\\
W_i\Omega_{i-1}-\Omega_iW_i=0,\ \Omega_{i-1}=0
&\Longrightarrow\Omega_i=0.
\end{align}
Thus $\mathcal N_W\succ0$. The prefix coordinates
$T_i=W_i\cdots W_1$ determine all gauge matrices by
$Q_i=T_i(T_i')^\dagger$. Hence
\begin{align}
\{D\pi_L[u]:u_2=\cdots=u_{L-1}=0\}
&\supseteq\{UZ_1+Z_2V^\top:Z_1\in\R^{r\times d},\
Z_2\in\R^{m\times r}\}\\
&=T_X\mathcal M_r .
\end{align}
Here the endpoint multipliers have column space $\qrange U$ and row
space $\qrange V^\top$, respectively; arbitrary endpoint perturbations
therefore realize both displayed summands. Surjectivity gives
\begin{align}
\dim\ker D\pi_L
&=rd+mr+(L-2)r^2-r(m+d-r)=(L-1)r^2\\
&=\dim\qrange\mathcal V_W .
\end{align}
The kernel inclusion and dimension identity prove equality.
Expanding $\mathcal V_W^*\mathcal V_W$ gives
Eq.~\eqref{eq:rigid-block-tridiagonal}.
\end{proof}

\subsection{Jointly balanced representatives}
\begin{theorem}[Canonical chain and residual orthogonal gauges]
For $X=U\Sigma V^\top$, define $D=\Sigma^{1/L}$ and
\begin{align}
W_1^0&=DV^\top,&W_i^0&=D\ (1<i<L),&
W_L^0&=UD .
\label{eq:rigid-canonical-chain}
\end{align}
Then
\begin{align}
\pi_L(W^0)&=X,&
(W_{i+1}^0)^\top W_{i+1}^0&=W_i^0(W_i^0)^\top=D^2,\\
\kappa(W_i^0)&=\kappa(X)^{1/L}.
\end{align}
All jointly balanced chains in a fixed regular product fiber differ by
$\mathcal G_L^O={\rm O}(r)^{L-1}$.
\end{theorem}
\begin{proof}
For a jointly balanced chain, put $H_i=W_iW_i^\top=W_{i+1}^\top W_{i+1}$.
Then
\begin{align}
H_iW_i&=W_iH_{i-1}\quad(1<i<L),\\
T_iT_i^\top&=H_i^i,\qquad T_i=W_i\cdots W_1,\\
XX^\top&=(W_LW_L^\top)^L,&
X^\top X&=(W_1^\top W_1)^L .
\end{align}
The identities follow inductively from
\begin{equation}
W_iH_{i-1}^{\,i-1}W_i^\top
=H_i^{\,i-1}W_iW_i^\top=H_i^i .
\end{equation}
If $W,W'$ have the same product, $T_i'=Q_i^{-1}T_i$, and both are balanced,
\begin{align}
W_1'^\top W_1'&=(X^\top X)^{1/L}=W_1^\top W_1,\\
W_1^\top(Q_1^{-\top}Q_1^{-1}-I)W_1&=0
\Longrightarrow Q_1Q_1^\top=I .
\end{align}
Using $W_i'=Q_i^{-1}W_iQ_{i-1}$ and $Q_{i-1}\in{\rm O}(r)$,
\begin{align}
W_i'^\top W_i'&=W_{i-1}'(W_{i-1}')^\top
=Q_{i-1}^\top W_i^\top W_iQ_{i-1}\\
&\Longrightarrow W_i^\top(Q_i^{-\top}Q_i^{-1}-I)W_i=0
\Longrightarrow Q_iQ_i^\top=I .
\end{align}
This proves the claim by induction.
\end{proof}

\subsection{Exact normal-equation spectrum and a fast spectral solve}
\begin{theorem}[Discrete-sine diagonalization]
\label{thm:joint-spectrum}
At Eq.~\eqref{eq:rigid-canonical-chain}, put
$d_a=\sigma_a^{1/L}$ and $\theta_k=k\pi/L$.
The complete spectrum of $\mathcal N_W$ is
\begin{equation}
\lambda_{abk}=d_a^2+d_b^2-2d_ad_b\cos\theta_k,
\quad a,b\leq r,\quad1\leq k<L.
\label{eq:rigid-chain-eigenvalues}
\end{equation}
Consequently
\begin{align}
\lambda_{\min}(\mathcal N_W)
&=2\sigma_r^{2/L}(1-\cos(\pi/L)),\\
\lambda_{\max}(\mathcal N_W)
&=2\sigma_1^{2/L}(1+\cos(\pi/L)),\\
\kappa(\mathcal N_W)
&=\kappa(X)^{2/L}\cot^2\!\left(\frac{\pi}{2L}\right).
\label{eq:rigid-chain-condition}
\end{align}
\end{theorem}
\begin{proof}
For each matrix entry $(a,b)$, write
$\omega_i=(\Omega_i)_{ab}$. Equation~\eqref{eq:rigid-block-tridiagonal} becomes
\begin{align}
(\mathcal N_W\omega)_i
&=(d_a^2+d_b^2)\omega_i-d_ad_b(\omega_{i-1}+\omega_{i+1}),\\
\omega_0&=\omega_L=0,\\
s_i^{(k)}&=\sqrt{2/L}\sin(ik\pi/L),\\
s_{i-1}^{(k)}+s_{i+1}^{(k)}
&=2\cos(k\pi/L)s_i^{(k)},&
\sum_{i=1}^{L-1}s_i^{(k)}s_i^{(\ell)}&=\delta_{k\ell}.
\end{align}
Hence $\mathcal N_Ws^{(k)}=\lambda_{abk}s^{(k)}$.
Moreover
\begin{align}
\lambda_{abk}
&=(d_a-d_b)^2+2d_ad_b(1-\cos\theta_k)\\
&\geq2\sigma_r^{2/L}(1-\cos(\pi/L)),\\
\lambda_{abk}
&\leq d_a^2+d_b^2+2d_ad_b\cos(\pi/L)\\
&\leq2\sigma_1^{2/L}(1+\cos(\pi/L)).
\end{align}
Equality is attained at $(a,b,k)=(r,r,1)$ and $(1,1,L-1)$.
\end{proof}

\begin{corollary}[Damped inverse and residual-to-projection error]
For $F_i=(\mathcal V_W^*u)_i$ and $\varepsilon\geq0$,
\begin{align}
\widehat F_{abk}
&=\sqrt{2/L}\sum_{i=1}^{L-1}\sin(ik\pi/L)(F_i)_{ab},\\
\widehat\Omega_{abk}
&=\frac{\widehat F_{abk}}{\lambda_{abk}+2\varepsilon},\\
(\Omega_i)_{ab}
&=\sqrt{2/L}\sum_{k=1}^{L-1}\sin(ik\pi/L)\widehat\Omega_{abk}.
\label{eq:rigid-dst-solve}
\end{align}
The canonical solve takes $O(Lr^2\log L)$ operations using sine transforms.
If $\|(\mathcal N_W+2\varepsilon I)\widetilde\Omega-F\|_F\leq\tau$,
\begin{align}
\nn{\widetilde\Omega-\Omega_\varepsilon}
&\leq\frac{\tau}{\lambda_{\min}+2\varepsilon},\\
\nn{\mathcal V_W(\widetilde\Omega-\Omega_\varepsilon)}
&\leq\frac{\sqrt{\lambda_{\max}}\tau}
{\lambda_{\min}+2\varepsilon},\\
\nn{\mathcal V_W(\Omega_\varepsilon-\Omega_0)}
&\leq\frac{2\varepsilon\sqrt{\lambda_{\max}}}
{\lambda_{\min}(\lambda_{\min}+2\varepsilon)}\nn F .
\end{align}
\end{corollary}
\begin{proof}
Use Eq.~\eqref{eq:rigid-chain-eigenvalues} entrywise and
\begin{equation}
\|\mathcal V_W\|^2=\lambda_{\max}(\mathcal N_W),\qquad
\left|\frac1{\lambda+2\varepsilon}-\frac1\lambda\right|
=\frac{2\varepsilon}{\lambda(\lambda+2\varepsilon)}.
\end{equation}
\end{proof}

\begin{corollary}[Depth asymptotics and CG complexity]
\begin{align}
1-\cos(\pi/L)&=\frac{\pi^2}{2L^2}+O(L^{-4}),\\
\lambda_{\min}(\mathcal N_W)
&=\frac{\pi^2\sigma_r^{2/L}}{L^2}(1+O(L^{-2})),\\
\kappa(\mathcal N_W)
&=\frac{4L^2}{\pi^2}\kappa(X)^{2/L}(1+O(L^{-2})).
\end{align}
For exact arithmetic conjugate gradients,
\begin{equation}
\|\Omega^{(k)}-\Omega_*\|_{\mathcal N}
\leq2\left(\frac{\sqrt{\kappa(\mathcal N)}-1}
{\sqrt{\kappa(\mathcal N)}+1}\right)^k
\|\Omega^{(0)}-\Omega_*\|_{\mathcal N}.
\end{equation}
These are normal-solve conditioning statements, not task-speedup guarantees.
\end{corollary}

\subsection{Product metric and optimization spectrum of a depth-$L$ chain}
\begin{theorem}[Depth-dependent product preconditioner]
\label{thm:chain-product-metric}
At a jointly balanced representative,
\begin{align}
K_X^{(L)}(G)
&=\sum_{i=1}^L L_iL_i^\top G R_i^\top R_i\\
&=\sum_{i=1}^L(XX^\top)^{(L-i)/L}
G(X^\top X)^{(i-1)/L},\\
g_X^{(L)}(Z,Y)&=\ii Z{(K_X^{(L)})^\dagger Y}.
\end{align}
Zero exponents in this formula denote $I_m$ or $I_d$.
The tangent eigenvalues are
\begin{align}
\chi_{ab}^{(L)}
&=\sum_{i=1}^L\sigma_a^{2(L-i)/L}\sigma_b^{2(i-1)/L},\\
\chi_{a,\perp}^{(L)}&=\sigma_a^{\,2-2/L},&
\chi_{\perp,b}^{(L)}&=\sigma_b^{\,2-2/L},\\
\chi_{ab}^{(L)}
&=\begin{cases}
\displaystyle\frac{\sigma_a^2-\sigma_b^2}
{\sigma_a^{2/L}-\sigma_b^{2/L}},&\sigma_a\ne\sigma_b,\\
L\sigma_a^{\,2-2/L},&\sigma_a=\sigma_b.
\end{cases}
\label{eq:rigid-chain-metric-spectrum}
\end{align}
In particular,
\begin{equation}
a^{2-2/L}\|P_XG\|_F^2
\leq\|\qgrad_{g^{(L)}}f\|_{g^{(L)}}^2
\leq Lb^{2-2/L}\|P_XG\|_F^2 .
\end{equation}
\end{theorem}
\begin{proof}
For $J_W=D\pi_L(W)$,
\begin{align}
(J_W^*G)_i&=L_i^\top G R_i^\top,&
J_WJ_W^*G&=\sum_i L_iL_i^\top G R_i^\top R_i .
\end{align}
At the canonical chain, the top-left block acts entrywise by
$\chi_{ab}^{(L)}$; the two off-diagonal tangent blocks receive only the
first and last summands, respectively. Finally,
\begin{equation}
\sum_{j=0}^{L-1}u^{L-1-j}v^j
=\frac{u^L-v^L}{u-v},\qquad
\lim_{v\to u}\frac{u^L-v^L}{u-v}=Lu^{L-1}.
\end{equation}
\end{proof}

\subsection{Aligned spectral dynamics and a finite-horizon obstruction}
\begin{proposition}[Exact aligned dynamics]
Assume fixed singular bases and
$G=U\qdiag(\gamma_i)V^\top$. Then
\begin{align}
\dot X&=-K_X^{(L)}G,&
\dot\sigma_i&=-L\sigma_i^{\,2-2/L}\gamma_i.
\end{align}
For $f(X)=\tfrac12\|X-X_*\|_F^2$ with
$X_*=U\qdiag(s_i)V^\top$,
\begin{align}
\dot\sigma_i&=L\sigma_i^{\,2-2/L}(s_i-\sigma_i),\\
\sigma_i^+
&=\sigma_i\left[1+\eta(s_i-\sigma_i)
\sigma_i^{(L-2)/L}\right]^L
\label{eq:rigid-aligned-discrete}
\end{align}
provided the factor update preserves positive singular coordinates.
\end{proposition}
\begin{proof}
For diagonal balanced factors $w_i=\sigma_i^{1/L}$,
\begin{align}
\partial_{w_i}f&=(\sigma_i-s_i)\sigma_i^{(L-1)/L},\\
w_i^+&=\sigma_i^{1/L}
-\eta(\sigma_i-s_i)\sigma_i^{(L-1)/L},&
\sigma_i^+&=(w_i^+)^L.
\end{align}
\end{proof}

\begin{corollary}[Small-mode hitting time]
Let $0<\sigma_0<\theta s<s$, $0<\theta<1$.
For $L=2$,
\begin{equation}
\sigma(t)=\frac{s}{1+(s/\sigma_0-1)e^{-2st}} .
\end{equation}
For $L>2$, $p=2-2/L>1$, the time $\tau_\theta$ to reach $\theta s$ satisfies
\begin{align}
\tau_\theta&=\int_{\sigma_0}^{\theta s}\frac{du}{Lu^p(s-u)},\\
\frac{\sigma_0^{1-p}-(\theta s)^{1-p}}{Ls(p-1)}
&\leq\tau_\theta\leq
\frac{\sigma_0^{1-p}-(\theta s)^{1-p}}{L(1-\theta)s(p-1)} .
\label{eq:rigid-small-mode-time}
\end{align}
\end{corollary}
\begin{proof}
On $u\in[\sigma_0,\theta s]$,
$(1-\theta)s\leq s-u\leq s$. Integrate
$u^{-p}/(Ls)$ and $u^{-p}/(L(1-\theta)s)$.
\end{proof}

\subsection{Global orbit energy of a deep factorization}
\begin{theorem}[Schatten variational identity and factor conditioning]
Let $p=2/L$. Over all compatible full-rank chains with product $X$,
\begin{align}
\inf_{\pi_L(W)=X}\frac1L\sum_{i=1}^L\nn{W_i}^2
&=\sum_{a=1}^r\sigma_a(X)^{2/L}
=\|X\|_{S_p}^{\,p},\\
\inf_{\pi_L(W)=X}\max_i\kappa(W_i)
&=\kappa(X)^{1/L}.
\label{eq:rigid-chain-orbit-energy}
\end{align}
The canonical jointly balanced chain attains both infima.
\end{theorem}
\begin{proof}
Multiplicative singular-value majorization and scalar H\"older imply
\begin{align}
\sum_a\sigma_a(X)^{2/L}
&\leq\sum_a\prod_{i=1}^L\sigma_a(W_i)^{2/L}\\
&\leq\prod_{i=1}^L
\left(\sum_a\sigma_a(W_i)^2\right)^{1/L}\\
&=\prod_{i=1}^L\nn{W_i}^{2/L}
\leq\frac1L\sum_i\nn{W_i}^2,\\
\kappa(X)&\leq\prod_i\kappa(W_i)
\leq\left(\max_i\kappa(W_i)\right)^L.
\end{align}
At Eq.~\eqref{eq:rigid-canonical-chain},
\begin{equation}
\nn{W_i^0}^2=\sum_a\sigma_a^{2/L},
\qquad\kappa(W_i^0)=\kappa(X)^{1/L}.
\end{equation}
\end{proof}
